\begin{helper}
Assume the hypotheses of \(\noderef{ProjectionClosedPairLossBound}\).  Let
\[
P_I=\noderef{TwoBiteProjectionPairSet}(\omega,I),\qquad
C_I=\{e\in P_I:\noderef{TwoBiteProjectionPairClosed}(\omega,I,e)\},
\]
put \(k=|I|\), and write
\[
\ell_R=|\pi_R(I)|,\qquad \ell_B=|\pi_B(I)|.
\]
Then
\[
|C_I|
\le
|P_I|-
\noderef{ProjectionOpenPairFunction}(\ell_R,\ell_B,k,p,n)
+9\varepsilon_1 k^2 .
\]
\end{helper}
\begin{proof}
Let \(f=\noderef{ProjectionOpenPairFunction}\).  By
\(\noderef{TwoBiteProjectionPairSet}\), the set \(P_I\) is the union of the
red-tagged base pairs in \(\pi_R(I)\) and the blue-tagged base pairs in
\(\pi_B(I)\).  The tags lie in different summands, so this union is disjoint.
Using \(\noderef{TwoBiteBasePairs}\), its cardinality is therefore
\[
|P_I|=\binom{\ell_R}{2}_{\mathbb R}
     +\binom{\ell_B}{2}_{\mathbb R}. \tag{1}
\]
Define
\[
A_R=
\min\left\{\binom{k-\ell_R}{2}_{\mathbb R},
  \binom{pn}{2}_{\mathbb R}+\binom{k-\ell_R-pn}{2}_{\mathbb R}\right\},
\]
\[
A_B=
\min\left\{\binom{k-\ell_B}{2}_{\mathbb R},
  \binom{pn}{2}_{\mathbb R}+\binom{k-\ell_B-pn}{2}_{\mathbb R}\right\}.
\]
Unfolding \(f\) and using (1) gives
\[
|P_I|-f(\ell_R,\ell_B,k,p,n)=A_R+A_B. \tag{2}
\]
It remains to show \(|C_I|\le A_R+A_B+9\varepsilon_1 k^2\).

Take any \(e\in C_I\).  By the definition of \(C_I\) and
\(\noderef{TwoBiteProjectionPairClosed}\), the pair \(e\) has a witnessing
base vertex \(x\).  Compare \(|X_x(I)|\) with the small, large, and huge
cutoffs.  Since \(X_x(I)\subseteq I\), the cardinality is between \(0\) and
\(k\).  Hence the witness is in exactly one of the four regimes represented
by \(\noderef{TwoBiteSmallClass}\), \(\noderef{TwoBiteMediumClass}\),
\(\noderef{TwoBiteLargeClass}\), or \(\noderef{TwoBiteHugeClass}\), allowing
overlap only at the level of different possible witnesses for the same
projected pair.  Thus the closed projected-pair set is contained in the union
of the small-, medium-, and large-witness projected-pair sets, the two
same-color huge projected-pair sets, and the two opposite-color huge
projected-pair sets.

For the large, medium, and small witness classes, apply
\(\noderef{ProjectedPairsInClassMultiplicityBound}\).  It bounds each
projected-pair set generated by one class by twice the corresponding
final-pair count:
\[
|Q_{\mathcal A}|
\le 2|\noderef{TwoBiteClosedPairsInClass}(\omega,I,\mathcal A)|
\quad
(\mathcal A=L_I,M_I,S_I).
\]
The finite union bound gives
\[
|Q_{L_I}\cup Q_{M_I}\cup Q_{S_I}|
\le 2A_L+2A_M+2A_S,
\]
where \(A_L,A_M,A_S\) are the three closed-pair counts for the large,
medium, and small classes.

Write
\[
H_I(x)\Longleftrightarrow \noderef{TwoBiteHugeClass}(\omega,I,x).
\]
The same-color huge red and blue pieces are also bounded by projected-pair
sums.  For example, let \(H_R^{=}\) be the red-tagged projected pairs in
\(C_I\) admitted by at least one huge red witness.  For each huge red witness
\(x\), the red-tagged pairs generated by this particular witness are exactly
the red tags of
\[
\noderef{TwoBiteBasePairs}
\bigl((\noderef{TwoBiteX}(\omega,I,x)).\mathrm{image}(\pi_R)\bigr).
\]
The set \(H_R^{=}\) is contained in the union of these finite sets as \(x\)
ranges over the huge red witnesses.  The finite union bound therefore gives
\[
|H_R^{=}|
\le
\sum_{\substack{x:\ H_I(x)\\ x\in V_R}}
  \binom{|\pi_R(X_x(I))|}{2}_{\mathbb R}.
\]
The same argument for huge blue witnesses gives the analogous bound for the
blue-tagged same-color piece by
\[
\sum_{\substack{x:\ H_I(x)\\ x\in V_B}}
  \binom{|\pi_B(X_x(I))|}{2}_{\mathbb R}.
\]
Unfolding \(\noderef{TwoBiteProjectedPairSum}\), these two sums are precisely
the huge-red-through-red and huge-blue-through-blue contributions.  Thus, if
\(A_H^{=}\) denotes their sum, then
\[
|C_{\mathrm{err}}|\le 2A_L+2A_M+2A_S+A_H^{=}.
\]
All four quantities are nonnegative, so the right side is at most
\[
2(A_L+A_M+A_S+A_H^{=}).
\]

The hypotheses needed by \(\noderef{ClosedPairsControlled}\) are now
available.  The large-class input is exactly
\(\noderef{LargeClosedPairsBound}\) applied with the present
\(\omega,I,\varepsilon,\varepsilon_1,\varepsilon_2,n_0\).  The medium and
small inputs are the assumed closed-pairs events at scale \(k=|I|\).  The
same-color huge input is the third conjunct of
\(\noderef{HugeProjectionClosedPairsBound}\), again applied with the present
parameters and the rounded identities for \(m,p,k\).  Therefore
\(\noderef{ClosedPairsControlled}\) gives
\[
A_L+A_M+A_S+A_H^{=}\le 4\varepsilon_1k^2.
\]
Combining this with the preceding display yields
\[
|C_{\mathrm{err}}|\le 8\varepsilon_1k^2. \tag{3}
\]

It remains to bound the two opposite-color huge pieces.  Define
\(C_R^{\mathrm{opp}}\) to be the blue-tagged projected pairs in \(C_I\) that
are admitted by at least one huge red witness, and define
\(C_B^{\mathrm{opp}}\) symmetrically as the red-tagged projected pairs
admitted by at least one huge blue witness.  The earlier witness
classification gives
\[
C_I\subseteq
C_{\mathrm{err}}\cup C_R^{\mathrm{opp}}\cup C_B^{\mathrm{opp}}. \tag{4}
\]
Fix a huge red witness \(x\).  The opposite-color pairs generated by this
particular witness are the blue tags of
\[
\noderef{TwoBiteBasePairs}
\bigl((\noderef{TwoBiteX}(\omega,I,x)).\mathrm{image}(\pi_B)\bigr).
\]
Therefore \(C_R^{\mathrm{opp}}\) is contained in the union of these tagged
base-pair sets over all \(x\) satisfying \(H_I(x)\) and \(x\in V_R\).  The
sets in this union need not be disjoint, but the finite union bound gives
\[
|C_R^{\mathrm{opp}}|
\le
\sum_{\substack{x:\ H_I(x)\\ x\in V_R}}
  \left|
  \noderef{TwoBiteBasePairs}
  \bigl((\noderef{TwoBiteX}(\omega,I,x)).\mathrm{image}(\pi_B)\bigr)
  \right|.
\]
For each fixed \(x\), the cardinality of the displayed base-pair set is
\[
\binom{|\pi_B(X_x(I))|}{2}_{\mathbb R}.
\]
Consequently
\[
|C_R^{\mathrm{opp}}|
\le
\sum_{\substack{x:\ H_I(x)\\ x\in V_R}}
  \binom{|\pi_B(X_x(I))|}{2}_{\mathbb R}. \tag{5}
\]
The predicate \(H_I(x)\land x\in V_R\) is exactly the Lean predicate
\[
\mathsf{hugeRed}(x)\Longleftrightarrow
\noderef{TwoBiteHugeClass}(\omega,I,x)\land
\exists r,\ x=\operatorname{inl}(r).
\]
Unfolding \(\noderef{TwoBiteProjectedPairSum}\), the right side of (5) is
therefore exactly
\[
\noderef{TwoBiteProjectedPairSum}(\omega,I,\mathsf{hugeRed},\pi_B).
\]
The fourth conjunct of \(\noderef{HugeProjectionClosedPairsBound}\) applies
to this sum and gives
\[
|C_R^{\mathrm{opp}}|\le (1+\varepsilon_1)A_R. \tag{6}
\]

The blue-witness case is identical with the colors reversed.  Namely,
\(C_B^{\mathrm{opp}}\) is contained in the union, over huge blue witnesses,
of the red-tagged base-pair sets generated by \(\pi_R(X_x(I))\).  The finite
union bound and the cardinality formula for
\(\noderef{TwoBiteBasePairs}\) give
\[
|C_B^{\mathrm{opp}}|
\le
\sum_{\substack{x:\ H_I(x)\\ x\in V_B}}
  \binom{|\pi_R(X_x(I))|}{2}_{\mathbb R}.
\]
The predicate \(H_I(x)\land x\in V_B\) is the Lean predicate
\[
\mathsf{hugeBlue}(x)\Longleftrightarrow
\noderef{TwoBiteHugeClass}(\omega,I,x)\land
\exists b,\ x=\operatorname{inr}(b).
\]
After unfolding \(\noderef{TwoBiteProjectedPairSum}\), the preceding sum is
\[
\noderef{TwoBiteProjectedPairSum}(\omega,I,\mathsf{hugeBlue},\pi_R).
\]
The fifth conjunct of \(\noderef{HugeProjectionClosedPairsBound}\) bounds
this sum, so
\[
|C_B^{\mathrm{opp}}|\le (1+\varepsilon_1)A_B. \tag{7}
\]
Together with (4), the finite union bound, (3), (6), and (7),
\[
|C_I|\le 8\varepsilon_1k^2+(1+\varepsilon_1)(A_R+A_B).
\]
Since \(\pi_R(I)\) and \(\pi_B(I)\) are images of the \(k\)-element set
\(I\), their deficits from \(k\) are nonnegative and at most \(k\).  The
first argument of each minimum defining \(A_R\) and \(A_B\) is therefore at
most \(\binom{k}{2}_{\mathbb R}\), and
\(\binom{k}{2}_{\mathbb R}\le k^2/2\).  Hence
\[
A_R+A_B\le k^2.
\]
Using \(\varepsilon_1\ge0\), this gives
\[
\varepsilon_1(A_R+A_B)\le \varepsilon_1k^2.
\]
Substituting in the previous estimate proves
\[
|C_I|\le A_R+A_B+9\varepsilon_1k^2.
\]
Finally replace \(A_R+A_B\) by \(|P_I|-f(\ell_R,\ell_B,k,p,n)\) using (2).
This is the asserted inequality.
\end{proof}
